% ========== EXTENSION PACKAGING & PUBLISHING ==========

\section{Extension Packaging \& Publishing}
\label{sec:extension-packaging-publishing}

\lettrine{T}{he} Extension Package Protocol (EPP) defines Symphony's packaging and publishing system, enabling secure, efficient distribution of extensions while maintaining compatibility, security, and performance across the ecosystem. The EPP format combines modern packaging techniques with Symphony-specific metadata and validation requirements.

\subsection{EPP Package Format}
\label{subsec:epp-package-format}

The Extension Package Protocol (EPP) format is designed as a secure, self-contained package that includes all necessary components for extension installation, validation, and execution.

\subsubsection{Package Structure}
\label{subsubsec:package-structure}

EPP packages follow a standardized directory structure that enables efficient processing and validation:

\begin{lstlisting}[caption=EPP Package Structure]
extension-name-1.2.3.epp
├── META-INF/
│   ├── MANIFEST.toml          # Extension manifest
│   ├── SIGNATURE.p7s          # Cryptographic signature
│   ├── CERTIFICATE.pem        # Signing certificate
│   └── CHECKSUMS.sha256       # File integrity checksums
├── src/                       # Source code (if applicable)
├── bin/                       # Compiled binaries
├── assets/                    # Static assets (icons, themes, etc.)
├── docs/                      # Documentation
├── tests/                     # Test suites
├── LICENSE                    # License text
└── README.md                  # Package documentation
\end{lstlisting}

\subsubsection{Package Metadata}
\label{subsubsec:package-metadata}

Each EPP package includes metadata that enables automated processing and validation:

\begin{table}[h]
\centering
\caption{EPP Package Metadata Components}
\label{tab:epp-metadata-components}
\begin{tabular}{@{}lll@{}}
\toprule
\textbf{Component} & \textbf{Purpose} & \textbf{Format} \\
\midrule
Manifest & Extension configuration & TOML \\
Signature & Cryptographic integrity & PKCS\#7 \\
Certificate & Identity verification & X.509 PEM \\
Checksums & File integrity & SHA-256 \\
Build Info & Compilation metadata & JSON \\
Dependencies & Dependency declarations & TOML \\
\bottomrule
\end{tabular}
\end{table}

\subsubsection{Binary Format Specification}
\label{subsubsec:binary-format-specification}

EPP packages use a ZIP-based container format with additional validation and security features:

\begin{infobox}[title=EPP Format Benefits]
\textbf{Compression:} Efficient storage and transfer using ZIP compression \\
\textbf{Integrity:} Built-in checksums and signature validation \\
\textbf{Security:} Cryptographic signatures and certificate chains \\
\textbf{Metadata:} Rich metadata for automated processing
\end{infobox}

\begin{lstlisting}[language=rust,caption=EPP Package Validation]
pub struct EPPPackage {
    manifest: ExtensionManifest,
    signature: CryptographicSignature,
    certificate_chain: Vec<Certificate>,
    content: PackageContent,
}

impl EPPPackage {
    pub fn validate(&self) -> ValidationResult {
        // 1. Verify package structure
        self.validate_structure()?;
        
        // 2. Validate manifest schema
        self.manifest.validate_schema()?;
        
        // 3. Verify cryptographic signature
        self.verify_signature()?;
        
        // 4. Check certificate chain
        self.validate_certificate_chain()?;
        
        // 5. Verify file checksums
        self.verify_checksums()?;
        
        // 6. Validate dependencies
        self.validate_dependencies()?;
        
        Ok(ValidationResult::Valid)
    }
}
\end{lstlisting}

\subsection{Publishing Workflow}
\label{subsec:publishing-workflow}

The publishing workflow guides extension developers through the process of preparing, validating, and distributing their extensions to the Symphony ecosystem.

\subsubsection{Development to Publication Pipeline}
\label{subsubsec:development-to-publication-pipeline}

\begin{figure}[htbp]
\centering
\begin{tikzpicture}[
    stage/.style={rectangle, rounded corners, draw=brandPrimary, fill=brandPrimary!10, text width=2.5cm, text centered, minimum height=1cm},
    validation/.style={diamond, draw=brandSecondary, fill=brandSecondary!10, text width=2cm, text centered, minimum height=0.8cm},
    flow/.style={->, >=stealth, thick, brandAccent}
]

% Pipeline stages
\node[stage] (develop) {Development};
\node[stage, right=2cm of develop] (build) {Build \& Test};
\node[stage, right=2cm of build] (package) {Package};
\node[validation, below=1.5cm of package] (validate) {Validate};
\node[stage, left=2cm of validate] (sign) {Sign};
\node[stage, left=2cm of sign] (publish) {Publish};

% Flow arrows
\draw[flow] (develop) -- (build);
\draw[flow] (build) -- (package);
\draw[flow] (package) -- (validate);
\draw[flow] (validate) -- node[right] {Pass} (sign);
\draw[flow] (sign) -- (publish);

% Failure loop
\draw[flow, dashed, red] (validate) to[bend right=45] node[below] {Fail} (develop);

% Stage details
\node[below=0.3cm of develop, text width=2.5cm, text centered, font=\small] {
    Code, test, \\
    document
};
\node[below=0.3cm of build, text width=2.5cm, text centered, font=\small] {
    Compile, \\
    run tests
};
\node[below=0.3cm of package, text width=2.5cm, text centered, font=\small] {
    Create EPP \\
    package
};

\end{tikzpicture}
\caption{Extension Publishing Pipeline}
\label{fig:publishing-pipeline}
\end{figure}

\subsubsection{Automated Build and Packaging}
\label{subsubsec:automated-build-packaging}

Symphony provides automated tools for building and packaging extensions:

\begin{lstlisting}[language=bash,caption=Symphony Extension Build Commands]
# Initialize new extension project
symphony-cli extension init --type instrument --name rust-generator

# Build extension for multiple platforms
symphony-cli extension build --target all --release

# Run test suite
symphony-cli extension test --coverage --integration

# Package extension into EPP format
symphony-cli extension package --sign --certificate dev-cert.pem

# Validate package before publishing
symphony-cli extension validate rust-generator-1.2.3.epp

# Publish to Symphony marketplace
symphony-cli extension publish --marketplace official
\end{lstlisting}

\subsubsection{Quality Assurance Gates}
\label{subsubsec:quality-assurance-gates}

The publishing workflow includes multiple quality assurance checkpoints:

\begin{table}[h]
\centering
\caption{Publishing Quality Gates}
\label{tab:publishing-quality-gates}
\begin{tabular}{@{}llll@{}}
\toprule
\textbf{Gate} & \textbf{Checks} & \textbf{Automated} & \textbf{Blocking} \\
\midrule
Code Quality & Linting, formatting, complexity & Yes & No \\
Security Scan & Vulnerability analysis, secrets & Yes & Yes \\
Test Coverage & Unit, integration, performance & Yes & Configurable \\
Manifest Validation & Schema, dependencies, permissions & Yes & Yes \\
Documentation & README, API docs, examples & Partial & No \\
License Compliance & License compatibility, attribution & Yes & Yes \\
\bottomrule
\end{tabular}
\end{table}

\subsection{Version Management}
\label{subsec:version-management}

Symphony's version management system supports the complex lifecycle of extensions while maintaining compatibility and enabling smooth upgrades.

\subsubsection{Version Lifecycle}
\label{subsubsec:version-lifecycle}

Extensions progress through multiple lifecycle stages, each with specific characteristics and policies:

\begin{figure}[htbp]
\centering
\begin{tikzpicture}[
    version/.style={rectangle, rounded corners, draw=brandSecondary, fill=brandSecondary!10, text width=2.5cm, text centered, minimum height=1cm},
    transition/.style={->, >=stealth, thick, brandPrimary}
]

% Version stages
\node[version] (dev) {Development \\ \small 1.0.0-dev};
\node[version, right=2.5cm of dev] (alpha) {Alpha \\ \small 1.0.0-alpha.1};
\node[version, right=2.5cm of alpha] (beta) {Beta \\ \small 1.0.0-beta.1};
\node[version, below=2cm of beta] (stable) {Stable \\ \small 1.0.0};
\node[version, left=2.5cm of stable] (maintenance) {Maintenance \\ \small 1.0.1};
\node[version, left=2.5cm of maintenance] (deprecated) {Deprecated \\ \small End-of-life};

% Transitions
\draw[transition] (dev) -- (alpha);
\draw[transition] (alpha) -- (beta);
\draw[transition] (beta) -- (stable);
\draw[transition] (stable) -- (maintenance);
\draw[transition] (maintenance) -- (deprecated);

% Parallel development
\draw[transition, bend left=30] (stable) to node[above] {Major} (dev);

\end{tikzpicture}
\caption{Extension Version Lifecycle}
\label{fig:version-lifecycle}
\end{figure}

\subsubsection{Release Channels}
\label{subsubsec:release-channels}

Symphony supports multiple release channels that cater to different user needs and risk tolerances:

\begin{table}[h]
\centering
\caption{Extension Release Channels}
\label{tab:release-channels}
\begin{tabular}{@{}llll@{}}
\toprule
\textbf{Channel} & \textbf{Stability} & \textbf{Update Frequency} & \textbf{Target Users} \\
\midrule
Stable & High & Monthly & Production users \\
Beta & Medium & Weekly & Early adopters \\
Alpha & Low & Daily & Testers, developers \\
Nightly & Experimental & Continuous & Extension developers \\
\bottomrule
\end{tabular}
\end{table}

\subsubsection{Backward Compatibility Management}
\label{subsubsec:backward-compatibility-management}

Symphony implements sophisticated backward compatibility management to minimize disruption during updates:

\begin{successbox}
\textbf{Compatibility Promise:} Extensions declare their compatibility guarantees, and Symphony enforces these promises through automated testing and validation.
\end{successbox}

\begin{compactlist}
\item \textbf{API Versioning}: Extensions can support multiple API versions simultaneously
\item \textbf{Deprecation Warnings}: Gradual deprecation with clear migration paths
\item \textbf{Compatibility Shims}: Automatic adaptation layers for minor incompatibilities
\item \textbf{Migration Tools}: Automated tools for updating configurations and data
\end{compactlist}

\subsection{Update Distribution}
\label{subsec:update-distribution}

The update distribution system ensures efficient, secure delivery of extension updates while minimizing disruption to user workflows.

\subsubsection{Distribution Architecture}
\label{subsubsec:distribution-architecture}

\begin{figure}[htbp]
\centering
\begin{tikzpicture}[
    component/.style={rectangle, rounded corners, draw=brandPrimary, fill=brandPrimary!10, text width=2.8cm, text centered, minimum height=1cm},
    cdn/.style={ellipse, draw=brandSecondary, fill=brandSecondary!10, text width=2cm, text centered, minimum height=0.8cm},
    flow/.style={->, >=stealth, thick, brandAccent}
]

% Distribution components
\node[component] (registry) {Extension Registry};
\node[cdn, right=3cm of registry] (cdn1) {CDN Node 1};
\node[cdn, above right=1cm and 1.5cm of cdn1] (cdn2) {CDN Node 2};
\node[cdn, below right=1cm and 1.5cm of cdn1] (cdn3) {CDN Node 3};
\node[component, right=3cm of cdn1] (client) {Symphony Client};

% Distribution flows
\draw[flow] (registry) -- (cdn1);
\draw[flow] (registry) to[bend left=15] (cdn2);
\draw[flow] (registry) to[bend right=15] (cdn3);
\draw[flow] (cdn1) -- (client);
\draw[flow] (cdn2) to[bend left=15] (client);
\draw[flow] (cdn3) to[bend right=15] (client);

% Labels
\node[above=0.3cm of registry] {\small Master Registry};
\node[above=0.3cm of cdn2] {\small Global CDN};
\node[above=0.3cm of client] {\small End Users};

\end{tikzpicture}
\caption{Extension Update Distribution Architecture}
\label{fig:update-distribution-architecture}
\end{figure}

\subsubsection{Delta Updates}
\label{subsubsec:delta-updates}

Symphony implements delta update technology to minimize bandwidth usage and update time:

\begin{infobox}[title=Delta Update Benefits]
\textbf{Bandwidth Efficiency:} Only changed files are downloaded \\
\textbf{Faster Updates:} Reduced download and installation time \\
\textbf{Reliability:} Checksums ensure integrity of delta patches \\
\textbf{Rollback Support:} Delta information enables quick rollbacks
\end{infobox}

\begin{table}[h]
\centering
\caption{Update Size Comparison}
\label{tab:update-size-comparison}
\begin{tabular}{@{}llll@{}}
\toprule
\textbf{Update Type} & \textbf{Typical Size} & \textbf{Download Time} & \textbf{Bandwidth Savings} \\
\midrule
Full Package & 50-200 MB & 30-120 seconds & 0\% (baseline) \\
Delta Update & 5-20 MB & 3-12 seconds & 80-90\% \\
Incremental Patch & 1-5 MB & 1-3 seconds & 95-98\% \\
\bottomrule
\end{tabular}
\end{table}

\subsubsection{Update Scheduling and Policies}
\label{subsubsec:update-scheduling-policies}

Symphony provides flexible update policies that balance security, stability, and user control:

\begin{lstlisting}[language=toml,caption=Update Policy Configuration]
[update_policy]
# Automatic update settings
auto_update_enabled = true
auto_update_channel = "stable"
auto_update_schedule = "daily"  # daily | weekly | monthly

# Update restrictions
update_during_work_hours = false
require_user_confirmation = ["major", "breaking"]
defer_updates_days = 7  # Maximum deferral period

# Rollback settings
enable_automatic_rollback = true
rollback_on_failure = true
rollback_timeout_minutes = 30

# Bandwidth management
use_metered_connection = false
max_download_speed_mbps = 10
update_only_on_wifi = true
\end{lstlisting}

\subsubsection{Update Verification and Rollback}
\label{subsubsec:update-verification-rollback}

Every update includes verification and rollback capabilities:

\begin{alertbox}
\textbf{Safe Updates:} All updates are verified before installation and can be automatically rolled back if issues are detected during the verification period.
\end{alertbox}

\begin{compactlist}
\item \textbf{Pre-Installation Verification}: Signature, compatibility, and dependency checks
\item \textbf{Post-Installation Validation}: Functionality tests and health checks
\item \textbf{Automatic Rollback}: Triggered by failures, crashes, or performance degradation
\item \textbf{Manual Rollback}: User-initiated rollback to previous versions
\item \textbf{Rollback Limits}: Configurable limits on rollback history retention
\end{compactlist}

The packaging and publishing system creates a robust foundation for Symphony's extension ecosystem, enabling secure, efficient distribution while maintaining the quality and reliability essential for professional development environments.